% Chapter 1
% 
\chapter{Introduction} % Main chapter title
\label{chap:Chapter1} % For referencing the chapter elsewhere, use Chapter~\ref{chap:Chapter1}


%-------------------------------------------------------------------------------
%---------
%
\section{High-Level Overview and Motivation} 
\label{sec:high_level_overview} 
The automotive industry is performing a profound transformation driven by the rise of the \gls{SDV}, 
autonomous driving technologies and increasingly interconnected \gls{ECU} \parencite{spglobalsdv}. 
With the grouth of these systems in complexity and responsability, the reliability of high-assurance software becomes critical. 
The systems must operate under strict real-time, reliability and safety requeriments regulated by standards as ISO 26262 \parencite{deloittesdvreadiness}. 
In this environment, memory safety has a special importance, as software faults can directly translate into physical hazards \parencite{adacore_nsa_cisa_msl}.

Historically, C++ has been the dominant language (paired with C) in embedded and automotive development due to its performance, 
deterministic resource use and mature ecosystem \parencite{ISO26262_Part6}. 
However, its dependence on manual memory management continues to be a major source of vulnerabilities such as: 
buffer overflows, dangling pointes, race conditions and undefined behaviors \parencite{MISRACpp2023}. 
Such issues are incompatible with the needs ofsafety critical cyber-physical systems, 
where a minor software defect can lead to serious consequences \parencite{AUTOSAR_Adaptive}.

Rust has rapidly emerged as a plausible successor for safety-critical software. 
Enforcing memory safety, ownerhip and data race free concurrency at compile time \parencite{docs_rs}. 
Rust promises to drastically reduce classes of defects edemic to C/C++ systems \parencite{ProgrammingRust}. 
Rust achieves this without introducing garbage collection, which can support predictable resource management in systems with real-time requirements \parencite{docs_rs}.

Despite the advantages Rust presents, the transition from C++ to Rust in large industrial ecosystems remains complex. 
Automotive software contains millions of lines of validated legacy C++ code. 
This code is oftenly built at top of complex frameworks and libraries as Boost C++ libraries \parencite{CXX_Crate}. 
Rewrinting those systems from scratch is infeasible. To avoid it, the companies usally pursue incremental and hybrid migration. 
This migration requires robust interoperability between Rust and C++ \parencite{Google_Rust_Interop}.

This dissertation is motivated byt the need to address that gap: 
enabling safe, practical and measurable integration of Rust into deeply entrenched C++ based automotive architectures, 
specially real-time related architectures.

Rewriting millions of lines of validated, legacy C++ code is operationally infeasible, 
necessitating a hybrid integration strategy. 
This transition is particularly problematic for automotive subsystems (such as sensor fusion, V2X communication and path planning) 
that are deeply coupled with the C++ ecosystem and reliant on complex, template-heavy frameworks like the Boost C++ Libraries.


\section{Problem Statement} 
\label{sec:problem_statement} %For referencing this section elsewhere, use Section~\ref{sec:problem_statement}

The central results from the tension between two different realities:
\begin{enumerate}
    \item \textbf{C++ is deeply entrenched in automotive systems}, specially through heavily templated, performance-critical frameworks like Boost.
    \item  \textbf{Rust provides strong memory safety guarantees}, but its adoption is hindered by the difficulty of interoperating with complex C++ abstrations creating a challenge for a transition between C++ and Rust. 
\end{enumerate} 

Currently, the interoperability approaches rely mainly on C-style \gls{FFI} . 
However, ist not possiblet to natively represent C++ templates, ownership semantics, 
RAII lifecycles or advanced Boost abstractions with C \gls{FFI} \parencite{CrubitDesign}.

As result of this limitation:
\begin{itemize}
    \item The integration often requires unsafe code, reducing Rust's safety benefits.
    \item Interfacing with templaty heavily libraries, as Boost.Geometry and Boost.Asio, introduces high engineering overhead.
    \item Crossing the Rust and C++ boundary imposes a performance cost that is poorly characterized for high-assurance systems.
    \item The reviewed literature does not provide a sufficiently detailed methodology for the systematic and incremental migration of C++ Boost components to Rust while maintaining compatibility with automotive subsystems.
\end{itemize}

The problem is not simply porting code from C++ to Rust, but the absence of validated techniques, architectural patterns, and tooling to bridge the impedance mismatch between C++ templates and Rust's ownership model, without breaking determinism, performance or safety requirements.


\section{Research Questions}
\label{sec:research_questions}

To address the research gap, this dissertation is guided by the following research questions:
\begin{enumerate}
    \item \textbf{RQ1:} What architectural and interoperability patterns can support the incremental migration of selected Boost components from C++ to Rust?
    \item \textbf{RQ2:} What criteria can be used to identify Boost components that are suitable candidates for incremental migration in automotive-oriented software?
    \item \textbf{RQ3:} How can selected Boost components be reimplemented using Rust abstractions while preserving their relevant observable behaviour?
    \item \textbf{RQ4:} How can Rust/C++ FFI boundaries be designed to explicitly manage ownership, representation, lifetime, and error-handling constraints?
    \item \textbf{RQ5:} What performance characteristics are observed when the selected Rust implementations are compared with their C++ counterparts under the evaluated workloads?
    \item \textbf{RQ6:} Which safety guarantees are provided by Rust within the migrated components, and which guarantees must remain explicit contracts at the FFI boundary?
\end{enumerate}

The questions are intentionally scoped to the selected case studies and evaluated workloads. In particular, the performance experiments measure observed execution time and throughput rather than worst-case execution time, and therefore do not establish hard real-time guarantees.

\section{Research Objectives}
\label{sec:research_objectives}

The primary objective of this dissertation is to design, implement and evaluate a methodology for the incremental migration of selected Boost C++ components into Rust, while maintaining controlled interoperability with existing C++ systems. The work focuses on identifying repeatable migration decisions, applying them to two representative low-level libraries, and evaluating the resulting technical trade-offs.

The specific objectives are:
\begin{itemize}
    \item Analyze existing interoperability mechanisms between Rust and C++, including C ABI based integration, \texttt{cxx}, \texttt{bindgen}, \texttt{autocxx}, and \texttt{cbindgen}.
    \item Define explicit criteria for selecting Boost components suitable for migration, including safety relevance, isolation, portability, interoperability, performance relevance, and implementation complexity.
    \item Select and reimplement representative functionality from \texttt{Boost.Align} and \texttt{Boost.Endian} using Rust-native abstractions rather than performing a purely syntactic translation.
    \item Design interoperability boundaries with explicit ownership, lifetime, representation, and error-handling contracts.
    \item Measure functional correctness, selected performance characteristics, memory behaviour, and FFI overhead for the implemented case studies.
    \item Identify migration patterns that are generalizable beyond the two selected libraries and distinguish them from decisions that are specific to their APIs or compatibility requirements.
    \item Characterize the limitations of the evaluated approach, particularly at unsafe Rust and Rust/C++ FFI boundaries.
\end{itemize}

\section{Research Contributions}
\label{sec:research_contributions}

The contributions of this dissertation are:
\begin{itemize}
    \item A structured set of criteria for assessing Boost components as incremental Rust migration candidates.
    \item A comparative analysis of Rust/C++ interoperability mechanisms with explicit attention to compatibility, safety, portability, performance, ABI stability, and build complexity.
    \item A migration methodology that emphasizes semantic analysis, Rust-native redesign, explicit interoperability contracts, validation, and critical analysis of residual risks.
    \item Rust-native reimplementations of selected functionality from \texttt{Boost.Align} and \texttt{Boost.Endian}, demonstrating two different FFI strategies.
    \item An analysis of unsafe-code and FFI boundaries, including the safety properties that remain enforceable inside Rust and those that become external contracts.
    \item A quantitative evaluation of the selected implementations, including execution time, throughput, memory observations, correctness checks, and FFI overhead.
\end{itemize}

\section{Dissertation Plan Structure}
\label{sec:dissertation_structure}

The remainder of this dissertation is organized as follows:
\begin{itemize}
    \item \textbf{Chapter 2 -- State of the Art:} Presents the technological background, related work, Rust/C++ interoperability mechanisms, Boost architecture, and the criteria used to compare interoperability approaches.
    \item \textbf{Chapter 3 -- Methodology and Solution Proposal:} Defines the component-selection criteria, migration-enabling rationale, and proposed architectural methodology.
    \item \textbf{Chapter 4 -- Migration Strategy:} Describes the systematic migration process used to analyse, redesign, implement, integrate, and validate selected Boost functionality.
    \item \textbf{Chapter 5 -- Selected Libraries:} Describes \texttt{Boost.Align} and \texttt{Boost.Endian}, their relevance to the study, and their selection as representative low-level case studies.
    \item \textbf{Chapter 6 -- Implementation, Evaluation and Critical Analysis:} Presents the Rust implementations, benchmark results, unsafe-code analysis, FFI safety analysis, and generalizable migration patterns.
    \item \textbf{Chapter 7 -- Conclusions and Future Work:} Summarizes the findings, answers the research questions, states the contributions and limitations, and identifies future work.
\end{itemize}

\section{Privacy and ethics}
In the context of this dissertation plan, which is focused in modernizing C++ based systems using Rust, 
ethical considerations were primarily centered on scientific integrity, software safety and intellectual property compliance.

The research relied on public open-source codebases, as Boost libraries.  No personal data, user-identifiable information (PII), or private operational logs were collected, processed, or stored. 
Consequently, the study falls outside the scope of GDPR requirements regarding human subjects.


\section{Use of AI Tools}
\label{sec:ai_tools}

For the preparation of this dissertation plan and supporting documentation, 
AI tools (specifically large language models) were used strictly for the following purposes:
\begin{itemize}
    \item \textbf{Document writing evaluation and improvement}, including linguistic refinement, clarity enhancement, and reorganization of text originally written by the author.
    \item \textbf{Assistance in identifying potential sources and directions for literature review}, such as suggesting relevant topics, keywords, and areas of research to investigate. All cited references included in the dissertation were individually verified, read, and selected by the author.
    \item Structural feedback and general writing suggestions for non-technical sections.
    \item Code example generation to support the investigation of the state of the art and also to facilitate the understanding.
\end{itemize}

The tools used were Perplexity (latest version), for deep search specialy for community based information, 
chat GPT (version GPT\-5.1) and Gemini (version 3), for document writing evaluation and improvement.

No AI-generated technical explanations, conceptual developments, design decisions, code, 
or analysis results were incorporated directly into the dissertation. 
All technical content is exclusively authored by the candidate.

The author takes full responsibility for the correctness, originality, and academic integrity of the work presented. If additional AI tools are used during the subsequent stages of the dissertation, this section will be updated accordingly to specify which tools were used, for what purpose, and in which components of the document. If no further AI usage occurs, this will also be explicitly stated.
